\definecolor{headerbg}{HTML}{d5f5e3}
\setlength{\arrayrulewidth}{1pt}

Les story points (SP) sont estimés selon la suite de Fibonacci. L'échelle retenue est : 1--2 (tâche simple), 3--5 (complexité modérée), 8--13 (tâche complexe ou incertaine). Ce sprint totalise \textbf{20 stories, 99 SP}.

\begin{longtable}{|>{\raggedright\arraybackslash}p{4.2cm}|>{\raggedright\arraybackslash}p{4.8cm}|>{\raggedright\arraybackslash}p{4.5cm}|c|}
\hline
\rowcolor{headerbg}
\textbf{User Story} & \textbf{Tâche} & \textbf{Critères d'acceptation} & \textbf{\rotatebox{90}{SP}} \\
\hline
\endfirsthead
\hline
\rowcolor{headerbg}
\textbf{User Story} & \textbf{Tâche} & \textbf{Critères d'acceptation} & \textbf{\rotatebox{90}{SP}} \\
\hline
\endhead
\hline
\endfoot

%------ US-021 ------
\multirow{2}{=}{En tant que Planificateur, je veux assigner manuellement une commande à un chauffeur et un véhicule afin d'organiser les livraisons.}
& Vérifier la disponibilité du chauffeur et la capacité du véhicule avant de valider l'affectation. & L'affectation est bloquée si la capacité du véhicule est dépassée ou si le chauffeur est indisponible. & 3 \\
\cline{2-4}
& Permettre au Planificateur de sélectionner un chauffeur et un véhicule pour une commande avec indication des conflits. & Le formulaire signale clairement les conflits et bloque la soumission en cas d'incompatibilité. & 2 \\
\hline

%------ US-022 ------
\multirow{2}{=}{En tant que Planificateur, je veux créer un plan de livraison quotidien regroupant plusieurs commandes afin d'organiser la journée.}
& Permettre au Planificateur de regrouper plusieurs commandes dans un plan journalier avec calcul des distances estimées et de la charge par chauffeur. & Le plan affiche le nombre de commandes, les distances estimées et la charge par chauffeur. & 3 \\
\cline{2-4}
& Permettre au Planificateur de réordonner visuellement les commandes dans le plan journalier. & La liste des commandes du plan est réordonnée par le Planificateur avec mise à jour immédiate. & 2 \\
\hline

%------ US-023 ------
\multirow{2}{=}{En tant que Planificateur, je veux voir la charge de chaque chauffeur afin d'équilibrer les affectations.}
& Calculer et afficher la charge de chaque chauffeur (nombre de commandes, distances, poids) pour la journée en cours. & Un avertissement est affiché si un chauffeur dépasse le seuil de charge configuré. & 3 \\
\cline{2-4}
& Afficher un indicateur de surcharge par chauffeur dans le tableau de bord du Planificateur. & L'indicateur change de couleur si le chauffeur est surchargé ; visible directement sur le tableau de bord. & 2 \\
\hline

%------ US-024 ------
\multirow{2}{=}{En tant que Planificateur, je veux vérifier la capacité d'un véhicule avant affectation afin d'éviter les surcharges.}
& Vérifier automatiquement la capacité maximale (poids et volume) du véhicule à chaque tentative d'affectation. & L'affectation est bloquée si le poids ou le volume dépasse la capacité maximale du véhicule. & 3 \\
\cline{2-4}
& Afficher au Planificateur les indicateurs de capacité restante lors du formulaire d'affectation. & Les indicateurs de capacité poids/volume restante s'affichent et se mettent à jour en temps réel. & 2 \\
\hline

%------ US-025 ------
\multirow{2}{=}{En tant que Planificateur, je veux réassigner une livraison en cas de problème afin d'assurer la continuité du service.}
& Permettre au Planificateur de réassigner une livraison à un autre chauffeur disponible, avec notification au nouveau chauffeur. & La réassignation met à jour le plan et notifie le nouveau chauffeur automatiquement. & 3 \\
\cline{2-4}
& Afficher au Planificateur la liste des chauffeurs disponibles pour une réassignation depuis la fiche livraison. & La fenêtre de réassignation s'affiche depuis la fiche livraison ; la liste des disponibles est filtrée. & 2 \\
\hline

%------ US-026 ------
\multirow{2}{=}{En tant que Planificateur, je veux que le système estime les distances entre les points de livraison afin d'optimiser les itinéraires.}
& Calculer les distances entre les points de livraison à partir des coordonnées géographiques enregistrées. & Les distances sont calculées et affichées ; la précision est acceptable pour les zones urbaines tunisiennes. & 5 \\
\cline{2-4}
& Valider les calculs de distance sur des jeux de données représentatifs de zones réelles. & Les distances calculées sont cohérentes avec les distances réelles sur les itinéraires utilisés. & 3 \\
\hline

%------ US-027 ------
\multirow{2}{=}{En tant que Planificateur, je veux gérer les créneaux horaires de livraison demandés par les clients afin de respecter leurs contraintes.}
& Intégrer les créneaux horaires demandés dans le modèle de commande et les vérifier lors de l'affectation. & Les créneaux sont respectés lors de l'affectation ; un conflit de créneau est signalé au Planificateur. & 2 \\
\cline{2-4}
& Permettre au REC de saisir un créneau horaire lors de la création ou modification d'une commande. & Le sélecteur de créneau est disponible dans le formulaire ; les conflits sont affichés avant soumission. & 1 \\
\hline

%------ US-028 ------
\multirow{2}{=}{En tant que Planificateur, je veux créer des plans de livraison hebdomadaires récurrents afin de gagner du temps.}
& Permettre au Planificateur de définir un modèle de plan hebdomadaire récurrent avec sélection des jours. & Le plan récurrent se génère automatiquement chaque semaine à partir du modèle défini. & 3 \\
\cline{2-4}
& Permettre au Planificateur d'activer ou désactiver un plan récurrent depuis son tableau de bord. & Le modèle de plan récurrent est sauvegardé et peut être activé ou désactivé. & 2 \\
\hline

%------ US-029 ------
\multirow{2}{=}{En tant que Planificateur, je veux générer un bon de chargement pour l'entrepôt afin de faciliter la préparation des expéditions.}
& Générer un document de chargement listant les articles par commande et par véhicule, sous format téléchargeable. & Le document liste les articles par commande et par véhicule ; téléchargeable en un clic. & 3 \\
\cline{2-4}
& Mettre à disposition du Planificateur le bouton de génération du bon de chargement depuis la page du plan journalier. & Le document est généré et téléchargeable depuis l'interface en un clic. & 2 \\
\hline

%------ US-030 ------
\multirow{2}{=}{En tant que Planificateur, je veux reprogrammer les livraisons échouées au jour suivant afin de finaliser les commandes.}
& Proposer automatiquement les livraisons échouées en priorité dans le plan du lendemain. & Les livraisons échouées sont proposées en priorité dans le plan du lendemain. & 3 \\
\cline{2-4}
& Afficher au Planificateur un panneau dédié aux livraisons à reprogrammer avec option de validation rapide. & Le panneau des livraisons à reprogrammer s'affiche avec option de validation rapide. & 2 \\
\hline

%------ US-031 ------
\multirow{2}{=}{En tant que Planificateur, je veux définir des règles d'incompatibilité entre chauffeurs et zones afin d'éviter les erreurs d'attribution.}
& Permettre au Planificateur de définir des règles d'incompatibilité entre un chauffeur et une zone, vérifiées à chaque affectation. & La règle est vérifiée lors de chaque affectation ; une alerte est levée en cas de violation. & 2 \\
\cline{2-4}
& Permettre au Planificateur de consulter et modifier la liste des règles d'incompatibilité définies. & La liste des règles est consultable et modifiable ; les violations sont signalées lors de l'affectation. & 1 \\
\hline

%------ US-032 ------
\multirow{2}{=}{En tant que Planificateur, je veux recevoir des suggestions d'itinéraires optimisés afin de minimiser les kilomètres parcourus.}
& Générer des itinéraires optimisés tenant compte de la capacité des véhicules, des créneaux horaires et des zones de livraison. & Le système produit des itinéraires optimisés ; le gain en kilomètres est calculé et affiché. & 8 \\
\cline{2-4}
& Afficher les suggestions d'itinéraires dans le tableau de bord du Planificateur avec un indicateur de qualité. & Les suggestions s'affichent dans le tableau de bord avec un score de qualité. & 5 \\
\hline

%------ US-033 ------
\multirow{2}{=}{En tant que Planificateur, je veux accepter, modifier ou rejeter chaque suggestion d'itinéraire afin de garder le contrôle sur les décisions.}
& Permettre au Planificateur de valider, ajuster ou rejeter chaque suggestion avec traçabilité pour améliorer les suggestions futures. & Le Planificateur peut valider, éditer ou rejeter ; l'action est enregistrée pour la rétroaction. & 3 \\
\cline{2-4}
& Mettre à disposition les actions (Accepter / Modifier / Rejeter) pour chaque suggestion avec possibilité d'ajustement. & Les actions sont accessibles sur chaque suggestion ; un formulaire d'ajustement est disponible. & 2 \\
\hline

%------ US-037 ------
\multirow{2}{=}{En tant que Planificateur, je veux recevoir des suggestions du chauffeur le plus adapté afin d'optimiser les affectations.}
& Calculer un score de pertinence pour chaque chauffeur selon ses performances, sa zone et sa disponibilité, et proposer les 3 meilleurs candidats. & Le système propose les 3 chauffeurs les plus adaptés avec un score de pertinence et une justification. & 3 \\
\cline{2-4}
& Afficher les suggestions de chauffeurs avec score et justification dans le formulaire d'affectation. & Les 3 suggestions s'affichent avec score et justification (zone, performance, disponibilité). & 2 \\
\hline

%------ US-034 ------
\multirow{2}{=}{En tant que Chauffeur, je veux voir l'itinéraire optimisé sur une carte afin de faciliter mes déplacements.}
& Afficher au Chauffeur les points de livraison numérotés dans l'ordre optimal sur une carte. & La carte affiche les points de livraison numérotés dans l'ordre optimal de passage. & 3 \\
\cline{2-4}
& Permettre au Chauffeur de lancer la navigation vers chaque point de livraison depuis l'application mobile. & Le Chauffeur peut lancer la navigation vers chaque point depuis l'application mobile. & 2 \\
\hline

%------ US-035 ------
\multirow{2}{=}{En tant que Planificateur, je veux que le système respecte toutes les contraintes métier lors de l'optimisation afin de garantir la faisabilité des suggestions.}
& Intégrer les contraintes métier dans le moteur d'optimisation (capacité, créneaux, zones, règles d'incompatibilité). & Aucune suggestion ne viole les contraintes définies ; un rapport de faisabilité est fourni. & 5 \\
\cline{2-4}
& Présenter au Planificateur un rapport des contraintes respectées et des éventuelles violations pour chaque suggestion. & Le rapport de faisabilité est consultable depuis l'interface de validation des suggestions. & 3 \\
\hline

%------ US-036 ------
\multirow{2}{=}{En tant que Planificateur, je veux que les parties prenantes soient notifiées après validation du plan afin d'assurer la coordination entre acteurs.}
& Envoyer une notification adaptée à chaque destinataire (Chauffeur, REC, Client) après validation du plan. & Le Chauffeur, le Responsable et le Client reçoivent chacun une notification adaptée à leur rôle. & 3 \\
\cline{2-4}
& Délivrer les notifications dans les 30 secondes suivant la validation du plan. & Les notifications apparaissent dans les 30 secondes suivant la validation. & 2 \\
\hline

%------ US-038 ------
\multirow{2}{=}{En tant que Planificateur, je veux envoyer des retours au système afin d'améliorer la qualité des suggestions futures.}
& Permettre au Planificateur d'enregistrer un retour (acceptation ou rejet avec motif) sur chaque suggestion du système. & Les retours sont enregistrés et utilisés pour améliorer les suggestions futures. & 3 \\
\cline{2-4}
& Afficher un formulaire de retour au Planificateur après chaque rejet d'une suggestion. & Le formulaire de motif est affiché après chaque rejet ; la saisie est optionnelle mais encouragée. & 2 \\
\hline

%------ US-106 ------
\multirow{2}{=}{En tant que Planificateur, je veux un tableau de bord opérationnel en temps réel afin de suivre l'avancement des livraisons.}
& Exposer au Planificateur les métriques opérationnelles en temps réel (livraisons en cours, retards, alertes). & Le tableau de bord s'actualise régulièrement ; les retards et alertes s'affichent automatiquement. & 3 \\
\cline{2-4}
& Afficher les indicateurs opérationnels (livraisons en cours, retards détectés, alertes) dans le tableau de bord. & Les indicateurs s'affichent avec mise à jour automatique. & 2 \\
\hline

%------ US-131 ------
\multirow{2}{=}{En tant que Chauffeur, je veux voir mes livraisons du jour et l'itinéraire optimisé sur la carte depuis l'application mobile.}
& Récupérer et afficher au Chauffeur la liste de ses livraisons du jour, triées dans l'ordre optimal défini par le Planificateur. & La liste des livraisons du jour est récupérée et triée dans l'ordre optimal défini par le Planificateur. & 3 \\
\cline{2-4}
& Afficher la carte avec les points de livraison numérotés et permettre la navigation vers chaque point. & La carte s'affiche avec les points numérotés ; la navigation est accessible pour chaque point. & 2 \\
\hline

\caption{Sprint Backlog du Sprint 3 --- Planification \& Optimisation (20 stories, 99 SP)}
\end{longtable}
\setlength{\arrayrulewidth}{0.4pt}
